\documentclass{article}
\usepackage{geometry}
\usepackage{concmath}
\usepackage[OT1]{fontenc}
\usepackage{graphicx}
\usepackage{amsfonts}
\usepackage{amsmath,mathtools}
\usepackage{calrsfs}
\usepackage{color}
\usepackage[dvipsnames]{xcolor}
\usepackage[many]{tcolorbox}
\usepackage{setspace}
\usepackage{tkz-tab}    
\usepackage{amssymb}          
\usepackage{multicol}  
\usepackage{esint}
\geometry{margin=0.5in}
\definecolor{green}{RGB}{52, 235, 168}
\definecolor{violet}{RGB}{165, 83, 237}
\definecolor{lemon}{RGB}{170, 237, 83}
\definecolor{red}{RGB}{240, 125, 84}
\definecolor{blue}{RGB}{105, 255, 242}



\begin{document}
    
\noindent \colorbox{green}{650.3}

* \(t \mapsto \frac{\sin t}{t}\) est continue sur \(]0, \infty[\)\\
* \underline{Au voisinnage de \(t \rightarrow 0\)}
\begin{align*}
    \frac{sint}{t} \backsim \frac{t}{t}\\
    & \rightarrow_{t \rightarrow 0} 1
\end{align*}
donc f se prolonge par continuité en 0 donc f est intégrable en 0 / \(\int_{0}^{1}f(t)dt\) est 
faussement généralisée \\

\noindent \underline{Au voisinnage de \(t \rightarrow +\infty\)}
Par parties, sous réserve de limites finies du crochet et de cv d'une intégrale 
\begin{align*}
  \int_{0}^{+\infty} \frac{1}{t} sint dt &= [\frac{1}{t}(1-cost)]^{+\infty}_{0} - \int_{0}^{\infty} (-) \frac{1}{t^{2}}(1-cost)dt \\ 
                                         &= 0 - 0 + \int_{0}^{\infty}\frac{1}{t^{2}}(1-cost)dt
\end{align*}

\noindent \colorbox{green}{650.4}
* \(f : x \mapsto ln(1 + \frac{sinx}{x})\) est continue sur \(]0, +\infty[\)\\
si \(x > 1, |\frac{sinx}{x}| < 1\) donc \(1 + \frac{\sin x}{x}>0\)\\
si \(x \leq 1, \frac{sinx}{x} \geq 0\)\\
* \underline{Au voisinnage de 0}\\
\begin{equation*}
    \frac{sinx}{x} \to 1
\end{equation*}
donc \(f(x) \to_{x \to 0} ln 2\)
donc f prolongeable par continuité en 0 donc intégrable en 0\\
* \underline{Au voisinnage de \(x \mapsto +\infty\)}\\
\begin{equation*}
    ln(1 + \frac{sinx}{x}) \sim \frac{sinx}{x}
\end{equation*}
méthode d'éclatement : DL(0,2) de f(x), f(x) converge \\

\noindent \colorbox{green}{650.28}\\
Au voisinnage de \(x \to 0\) \begin{equation*}
    f(x) = \int_{x}^{1}\frac{e^t}{dt}   
\end{equation*}
Intégrale partielle de \(\int_{\to 0}^{1}\frac{e^t}{dt}\)
Au voisinnage de \(t \rightarrow 0, \frac{e^{t}}{t} \backsim \frac{1}{t}\) Non intégrable en 0 donc \(f(x) \rightarrow_{x\to 0}+\infty\) par positivité \\ 
\begin{align*}
    f(x) = \int_{x}^{1} \frac{1}{t}dt + \int_{x}^{1} \frac{e^{t}-1}{t}dt \\ 
\end{align*}
intégrale partielle de \(t \rightarrow 0\) \(\frac{e^{t}-1}{t}\) = 1 + o(1) \\
donc \(g(x) \rightarrow l\)\\
\begin{align*}
  f(x) &= [lnt]^{1}_{x} + g(x)\\
       &= -lnx + g(x)\\
       &= -lnx + l + o(1)\\
       &\backsim_{x\to 0}-lnx 
\end{align*}
lnx est intégrable en 0 donc f est intégrable en 0 
\colorbox{blue}{Méthode 2}\\
\begin{align*}
  f(x) &= \int_{x}^{1} \frac{1}{t}e^{t} dt \\
       &= [lnt e^{t}]^{1}_{x} - \int_{x}^{1}lnt e^{t} dt \\
       &= -e^{x}lnx - \int_{x}^{1} lnt e^{t} dt \\
       &= -ln(x) + o(lnx)
\end{align*}






\end{document}

